\scanpage{164}
\mspage{19}
Ceci correspond à l'intuition (sortant immédiatement des exemples élémentaires), moyennant quoi, pour $k$ alg clos disons, les $H^0(K,F), H^1(K,F)\ldots$ ont une structure de groupes algébriques (ind-pro \ldots) sur $k$. Dans cette construction, le schéma en anneaux sur $k$ défini par $V$ (introduit par Serre) et le ``foncteur de Greenberg'' (associant à un schéma sur $V$ un proschéma sur $k$) joueront un rôle essentiel.

Ceci posé, la formule de dualité s'écrive formellement comme dans (D) ci-dessus
\[
 \Delta R\underline{\Gamma}_{K\unsure{.}}(F)\simeq R\underline{\Gamma}_K(DF[1])
\]
où $D$ est la dualité de Cartier, $\Delta$ est la dualité de Serre. Pour un corps résiduel fini, elle devient (via ``l'astuce de Lang'' plus haut)
\[
 \begin{aligned}
 \overset{\unsure{\sim}}{\Delta}_0 R\Gamma_K(F)\simeq R\Gamma_{K\unsure{.}}(DF[\deleted{\unsure{1}}2]),&\quad \Delta_0\ \text{dualité de}\\
 &\quad\text{\unsure{Pontrjagin|Pontryagin}.}
 \end{aligned}
\]
Elle contient encore le corps de classes local géométrique à la Serre, et le corps de classes local arithmetique\unsure{.} sous sa forme classique.

\uline{Remarques} \unsure{a|e|$\alpha$}) Si $F$ est premier à la car. résiduelle, la formule est très facile et bien connue. On peut la considérer comme un cas très particulier
% archive sheet number: 164
